\chapter{Conclusions and Future Work}
\label{ch:conclusions}

\section{Conclusions}
\label{sec:conclusions}

This thesis investigated whether structured graph knowledge and explicit empirical
statistical evidence can be integrated within a common modelling framework for two
complementary problems: reconstruction of directed graph relations and prediction of
patient-level clinical endpoints.

In Task~A, a heterogeneous graph encoder was combined with a candidate-specific
HCR-inspired statistical representation derived from patient observations. The results
showed that explicit empirical context provided substantial predictive information beyond
graph topology alone. The Last-FM* proxy demonstrated that the same broad modelling
principle could be transferred to a different domain and ranking objective, while the sealed
benchmark also showed that improvements within a neural graph model do not necessarily
translate into superiority over strong classical methods.

In Task~B, patient-specific observations were propagated over a shared pharmacotherapy
graph and complemented with endpoint-oriented statistical evidence. The results showed
that the value of this evidence depends on the ability of the graph neural network to
preserve predictive information across message-passing layers. Residual connections and
the statistical evidence branch can address related information-loss mechanisms on the
synthetic graph, while the Hetionet-derived proxy showed that their effects may become
complementary when relevant information must traverse longer or partially latent paths.

Taken together, the experiments support a methodological rather than clinical conclusion.
Structured graph representations and explicit empirical dependence can provide
complementary information, but their contribution depends on the graph topology,
prediction task, information pathway and evaluation protocol. The synthetic and proxy
experiments establish a foundation for further investigation on real clinical data; they do
not establish clinical utility or causal validity.

\section{Future Work}
\label{sec:future-work}

\paragraph{Real-world clinical validation.}
The most important next step is evaluation on real longitudinal clinical data. Such studies
would need to address irregular observation, missingness, heterogeneous coding,
treatment-selection effects, institutional variation and unmeasured confounding. The
primary question would be which components of the architecture remain useful when both
the graph structure and empirical evidence are uncertain.
Multicentre validation would be especially important for distinguishing relationships that
remain stable across institutions from correlations that reflect local observation or treatment
practice.

\paragraph{Higher-order statistical dependence.}
The present experiments make limited use of higher-order HCR coefficients, particularly
for continuous and mixed variables. Richer clinical datasets containing laboratory
measurements and longitudinal physiological variables would provide a more informative
setting in which to evaluate whether higher-order orthonormal coefficients capture
dependence patterns that cannot be represented adequately by first-order association
measures.

\paragraph{Causal discovery and graph refinement.}
A further direction is to combine mechanistic graph priors with dedicated causal-discovery
methods. HCR should remain an empirical dependence representation rather than a causal
estimator, but its coefficients could provide candidate statistical evidence to methods that
introduce the additional assumptions required for directional inference. Approaches based
on additive-noise models, LiNGAM-type methods, or other suitable causal-discovery
frameworks could therefore be investigated alongside graph priors, with careful attention to
variable type and assumption validity rather than automatic application across the graph.

\paragraph{Patient-specific mechanistic path activation.}
Clinical endpoint prediction alone does not identify which mechanisms are responsible
for the predicted risk. A natural extension is therefore to move from endpoint prediction
toward patient-specific activation of mechanistic paths.
For patient $p$ and endpoint $Y$, let
\begin{equation}
\mathcal{P}_p(Y)=\{\pi_1,\pi_2,\ldots,\pi_K\}
\end{equation}
denote plausible mechanistic routes, with patient-conditioned activation weights
$w_p(\pi_k)$.
Several routes should be allowed to remain simultaneously active. This is particularly
important in polypharmacy, where multiple drugs, risk factors and intermediate mechanisms
may contribute concurrently to the same adverse endpoint.
A future temporal aggregation layer could model how these routes accumulate, interact
and converge over time rather than forcing the system to select a single dominant pathway.

\paragraph{Active information acquisition and missing evidence.}
Missing information could also be treated as part of the reasoning problem rather than an
incidental preprocessing detail. When several mechanistic paths remain compatible with the
observed patient state, the model could identify which unobserved variable or clinical
measurement would be most informative for distinguishing between them.
Such an approach would connect prediction with targeted information acquisition:
the system would not only estimate risk, but identify which additional observation could
most reduce uncertainty about the mechanism responsible for that risk.

\paragraph{Entropy-based path weighting.}
Entropy-based graph-walk methods, including maximal-entropy random walks, provide one
possible starting point for weighting alternative routes through a mechanistic graph.
Standard topology-driven random walks would, however, require modification for the
present setting because clinical path relevance should depend on patient-specific evidence,
not graph structure alone.
A future patient-conditioned finite-path formulation could therefore combine structural
accessibility with observed activation of the variables along a route. Such a method should
permit several routes to contribute simultaneously to the same endpoint.

\paragraph{KAN as an alternative nonlinear encoder.}
Kolmogorov--Arnold Networks (KANs) could be evaluated more systematically as an
alternative to the MLP blocks used for statistical encoding and fusion. Their learnable
univariate transformations may be attractive for heterogeneous statistical descriptors and
nonlinear interactions.
Exploratory MLP-versus-KAN comparisons for Task~A role encoders did not establish an
advantage of KAN over MLPs under the reported protocol, so this remains an open
architectural comparison rather than a settled design choice.

\paragraph{Towards mechanistic digital twins.}
The combination of patient-specific graph instantiation, mechanistic prior knowledge and
explicit empirical evidence is conceptually related to emerging medical digital-twin
research. The goal should not be understood as constructing a complete virtual replica of
a human patient. A more realistic intermediate objective is a dynamically updated
representation of the clinically relevant mechanisms, exposures, observations and
uncertainties associated with a specific decision problem.
A clinically useful digital-twin framework would ultimately require temporal updating,
intervention modelling, uncertainty quantification, external validation and careful
governance. The present thesis addresses only some of the representational prerequisites
for such a system.

\paragraph{Iterative Task~A--Task~B loop.}
Finally, Task~A and Task~B could be connected into an iterative framework. High-confidence
candidate relations identified during graph reconstruction could be introduced into a
provisional graph and subsequently evaluated through patient-level endpoint prediction.
Conversely, systematic endpoint-prediction errors could identify regions of the graph in
which structural knowledge requires further examination.
Such a loop would move the framework from static graph reconstruction and endpoint
prediction toward continual evaluation of mechanistic knowledge against accumulating
empirical evidence.

The longer-term objective is therefore not only to predict what may happen, but to connect
prediction with the mechanisms, uncertainty and potentially modifiable pathways that make
the prediction clinically meaningful.
